\section{A uniform Gaussian limit for the actual return record}
\label{sec:stopped-gaussian}
We first obtain a quantitative characteristic-function estimate on the
collision clock. Smooth approximation is used only inside the proof;
the limiting statement concerns the original record and the original
section probability.

\subsection{Smooth perturbation with a shrinking scale}
\begin{lemma}[A quantitative smooth collision expansion]
\label{lem:smooth-collision-expansion}
On each of the collision spaces of
Lemma~\ref{lem:collision-spectral-input}, set
\[
 \mathcal L_{R,\delta,z}
       =\mathcal L_R M_{\exp(i z\cdot h_{R,\delta})},
       \qquad z\in\mathbb C^4.
\]
There are $a,C>0$ and $\rho<1$, independent of $R$ and
$0<\delta<1/2$, such that for $|z|<a\delta^2$ this family has a simple
analytic eigenvalue $\lambda_{R,\delta}(z)$ near one and an analytic
rank-one projection $\Pi_{R,\delta}(z)$. Its complementary powers are
bounded by $C\rho^m$, and, on a smaller such ball,
\begin{equation}\label{eq:smooth-cubic-expansion}
 \begin{split}
 \log\lambda_{R,\delta}(z)
       &=-\tfrac12z^{\mathsf T}\Gamma_{R,\delta}z
                  +O(\delta^{-6}|z|^3),\\
 \|\Pi_{R,\delta}(z)-\Pi_R\|
       &\le C\delta^{-2}|z| .
 \end{split}
\end{equation}
All constants are uniform. The logarithm is the branch vanishing at
$z=0$.
\end{lemma}
\begin{proof}
The smooth multiplier estimate and the smoothing bounds give
$\|\mathcal L_{R,\delta,z}-\mathcal L_R\|
\le C|z|\delta^{-2}$ when $|z|\delta^{-2}$ is small.
The uniform resolvent contours around one and around the complementary
spectrum in the proof of Proposition~\ref{prop:soft-killing} therefore
apply on $|z|<a\delta^2$. The resolvent Neumann series gives the analytic
projection, its difference estimate, and complementary powers on a
fixed circle of radius $\rho<1$. Decrease $a$ so that the eigenvalue
stays in a fixed disk about one not containing zero.

The eigenvalue derivative at zero is zero since $\int h_{R,\delta}
\dd\nu=0$. We identify its second derivative, rather than assuming a
variance formula. Fix a real direction $v$, put $u=v\cdot h_{R,\delta}$,
and normalize a right eigenvector $e(t)$ by $\ell_j(e(t))=1$.
Differentiating the eigenvector equation at zero yields
\[
 e'(0)=i\sum_{k\ge1}\mathcal L_R^k(u\nu),
\]
where the series converges in $\mathcal B_j$ by the spectral splitting.
A second differentiation followed by the mass functional gives
\[
 \lambda''(0)=-\int u^2\dd\nu
       -2\sum_{k\ge1}\int u(u\circ T_R^k)\dd\nu
       =-v^{\mathsf T}\Gamma_{R,\delta}v.
\]
Polarization identifies the full Hessian. Since the first derivative
vanishes, the Hessian of the logarithm is the same. Cauchy's estimates
for the bounded analytic logarithm on a ball of radius $a\delta^2$
bound its third derivatives by $C\delta^{-6}$ on a smaller ball.
Taylor's formula proves \eqref{eq:smooth-cubic-expansion}.
\end{proof}

\begin{proposition}[Uniform collision characteristic estimate]
\label{prop:collision-gaussian}
For every $K<\infty$ there is $C_K$ such that, for $m\ge2$ and for
$\sigma_R=1$ or $\sigma_R=s_R$,
\begin{equation}\label{eq:collision-gaussian}
 \sup_{R\in I,\ |t|\le K}
 \left|\int \sigma_R e^{it\cdot S_{m,R}h_R/\sqrt m}\dd\nu
       -e^{-t^{\mathsf T}\Gamma_Rt/2}\right|
       \le C_K m^{-1/26}\sqrt{\log(2+m)}.
\end{equation}
\end{proposition}
\begin{proof}
Smooth both $h_R$ and the initial density at scale $\delta$. The
smoothed density $\sigma_{R,\delta}=\mathsf S_\delta\sigma_R$ is
nonnegative, has integral one, and satisfies
$\|\sigma_{R,\delta}\nu\|_{\mathcal B_j}\le C\delta^{-2}$.
For $z=t/\sqrt m$, the transfer pairing and the spectral decomposition
of Lemma~\ref{lem:smooth-collision-expansion} give
\[
 \int\sigma_{R,\delta}e^{it\cdot S_{m,R}h_{R,\delta}/\sqrt m}\dd\nu
 =\lambda_{R,\delta}(z)^m
       \ell_j\bigl(\Pi_{R,\delta}(z)\sigma_{R,\delta}\nu\bigr)
       +O(\delta^{-2}\rho^m).
\]
At $z=0$ the amplitude is one; its difference from one is
$O_K(\delta^{-4}m^{-1/2})$. The logarithmic expansion gives, whenever
$K/\sqrt m<a\delta^2$ and $\delta^{-6}m^{-1/2}$ is sufficiently small,
\begin{equation}\label{eq:smoothed-characteristic-error}
 \left|\int\sigma_{R,\delta}
 e^{it\cdot S_{m,R}h_{R,\delta}/\sqrt m}\dd\nu
       -e^{-t^{\mathsf T}\Gamma_{R,\delta}t/2}\right|
 \le C_K\bigl(\delta^{-6}m^{-1/2}
          +\delta^{-4}m^{-1/2}+\delta^{-2}\rho^m\bigr).
\end{equation}
Here positive semidefiniteness of $\Gamma_{R,\delta}$ bounds the
Gaussian factor by one. The estimate for the exponential of the
Taylor remainder follows, for example, from
$|e^w-1|\le |w|e^{|w|}$.

The change of initial density costs at most $C\delta$ in the original
pairing, because the exponential has modulus one. The centered
residual $r_{R,\delta}=h_R-h_{R,\delta}$ has uniformly bounded
supremum and variation norms and $L^1$ norm $O(\delta)$. Therefore
Proposition~\ref{prop:bv-correlation}, component by component, gives
\[
 \int\sigma_{R,\delta}
       |S_{m,R}r_{R,\delta}|^2\dd\nu
       \le C m\delta(1+|\log\delta|),
\]
since the initial density is bounded independently of $\delta$.
The inequality $|e^{ix}-e^{iy}|\le|x-y|$ and Cauchy--Schwarz bound
the change of observable by
$C_K\sqrt{\delta(1+|\log\delta|)}$. Finally,
Theorem~\ref{thm:collision-covariance} bounds the change of the Gaussian
factor by $C_K\delta(1+|\log\delta|)$.

Take $\delta=\tfrac14 m^{-1/13}$. The smoothing error is
$O(m^{-1/26}\sqrt{\log(2+m)})$ and the cubic spectral error is
$O(m^{-1/26})$. The two smallness requirements hold for all sufficiently
large $m$, depending only on $K$. The amplitude and complementary errors
are smaller. Increasing $C_K$ handles the remaining finitely many $m$.
This proves the estimate uniformly over the finite cover of collision
spaces. It does not use an induced spectral gap.
\end{proof}

\subsection{Stopping at the genuine return times}
\begin{lemma}[Exact compensation and a return-clock window]
\label{lem:stopped-compensation}
For $\nu_R^*$-almost every initial state and every $n\ge1$,
\begin{equation}\label{eq:exact-stopped-compensation}
 J_{n,R}-n\bar G_R=S_{N_{n,R},R}h_R.
\end{equation}
There is $C$ such that for $n\ge2$ and $2\le a\le n$,
\begin{equation}\label{eq:return-clock-window}
 \nu_R^*\{|N_{n,R}-n/c_*|>a\}\le Cn/a^2.
\end{equation}
\end{lemma}
\begin{proof}
A trajectory starting in $Y_R^*$ has exactly $n$ visits to this section
among collision times $0,\ldots,N_{n,R}-1$. These are its initial visit
and its first $n-1$ returns. Also
$J_{n,R}=S_{N_{n,R},R}f_R$. Substitution in
\eqref{eq:bounded-compensation} proves the identity with no remainder.

Let $H_{L,R}=\sum_{j=1}^L\eta_R\circ T_R^j$. Its stationary mean is
$c_*L$. Proposition~\ref{prop:bv-correlation} applied to
$\eta_R-c_*$ gives
\[
 \int|H_{L,R}-c_*L|^2\dd\nu\le CL.
\]
The same upper bound, with a changed constant, holds under $\nu_R^*$
because its density relative to $\nu$ is at most $c_*^{-1}$.
The exact equivalences
$N_{n,R}>L\iff H_{L,R}<n$ and
$N_{n,R}\le L\iff H_{L,R}\ge n$ reduce both tails to this estimate.
Choose integers within distance one of $n/c_*+a$ and $n/c_*-a$.
The corresponding deviations from $c_*L$ are at least $c_*a-O(1)$,
and $L\le Cn$. Chebyshev's inequality proves the assertion when $a$
exceeds a fixed constant. For the bounded remaining range, enlarge
$C$ so that the right side is at least one.
\end{proof}

\begin{lemma}[A deterministic-window maximal estimate]
\label{lem:dyadic-collision-maximal}
For $L\ge1$, $r\ge0$, and $R\in I$,
\begin{equation}\label{eq:dyadic-collision-maximal}
 \int\max_{0\le j\le L}
 |S_{r+j,R}h_R-S_{r,R}h_R|^2\dd\nu
       \le CL[\log(2+L)]^2.
\end{equation}
The same estimate holds under $\nu_R^*$ and for backward sums on a
deterministic interval contained in the nonnegative collision times.
\end{lemma}
\begin{proof}
Absolute summability of the correlations gives second moment at most
$C\ell$ for the sum over every interval of $\ell$ collisions, by
stationarity. Enlarge $L$ to a power of two smaller than $2L$.
Every prefix of that interval is a disjoint union of at most one dyadic
interval at each scale. Cauchy--Schwarz bounds the square of its sum by
the number of scales times the sum of the squares of these interval
sums. Bound the latter by the sum over all dyadic intervals at all
scales, independently of the prefix. At each scale, the expectations
sum to at most $CL$, since the intervals partition the full interval.
There are $O(\log(2+L))$ scales, proving the bound. The proof applies
to each component and to intervals in reverse order. A bounded initial
density gives the assertion under $\nu_R^*$ as well.
\end{proof}

\begin{theorem}[Parameter-uniform Gaussian limit of the physical record]
\label{thm:actual-record-gaussian}
Define $D_R=c_*^{-1}\Gamma_R$. It is a continuous positive semidefinite
matrix on $I$. For every $K<\infty$ there is $C_K$ such that
\begin{equation}\label{eq:actual-record-gaussian}
 \sup_{R\in I,\ |t|\le K}
 \left|\int_{Y_R^*}
 e^{it\cdot(J_{n,R}-n\bar G_R)/\sqrt n}\dd\nu_R^*
          -e^{-t^{\mathsf T}D_Rt/2}\right|
       \le C_K n^{-1/26}\sqrt{\log(2+n)}
       \qquad(n\ge2).
\end{equation}
In particular the original, unsmoothed return record has a centered
Gaussian weak limit with covariance $D_R$. The estimate is uniform in
the physical parameter, including along convergent parameter sequences.
\end{theorem}
\begin{proof}
Put $m_n=\lfloor n/c_*\rfloor$ and
$a_n=\lceil n^{3/5}\rceil$. The clock-window lemma, with a harmless
change for rounding, bounds the probability of
$|N_{n,R}-m_n|>a_n$ by $Cn/a_n^2$. On the complementary event, the
difference $S_{N_{n,R},R}h_R-S_{m_n,R}h_R$ is bounded by the larger
of the forward and backward maxima on intervals of length $a_n$ about
$m_n$. The intervals have nonnegative endpoints since $c_*<1$.
The preceding lemma and Cauchy--Schwarz consequently give
\begin{equation}\label{eq:stopped-characteristic-error}
 \begin{split}
 &\left|\int e^{it\cdot S_{N_{n,R},R}h_R/\sqrt n}\dd\nu_R^*
       -\int e^{it\cdot S_{m_n,R}h_R/\sqrt n}\dd\nu_R^*\right|\\
 &\hspace{20mm}\le C_K\left(\frac{n}{a_n^2}
                  +\sqrt{\frac{a_n}{n}}\log(2+a_n)\right)
       \le C_K n^{-1/5}\log(2+n).
 \end{split}
\end{equation}
The exceptional event costs at most twice its probability. No estimate
of an unbounded sum on that event is used.

Apply Proposition~\ref{prop:collision-gaussian} at time $m_n$ and
frequency $t\sqrt{m_n/n}$ with initial density $s_R$. This frequency
stays in a fixed compact set. Since $m_n/n=c_*^{-1}+O(n^{-1})$ and
$\Gamma_R$ is uniformly bounded, its Gaussian differs from the one in
\eqref{eq:actual-record-gaussian} by $O_K(n^{-1})$. The identity
\eqref{eq:exact-stopped-compensation} completes the estimate; the error
in \eqref{eq:stopped-characteristic-error} is smaller than the stated
one. The continuity of $D_R$ follows from
Theorem~\ref{thm:collision-covariance}. The characteristic-function
continuity theorem yields the weak limit, allowing a singular Gaussian.
Along $R_n\to R_0$, the same bound and continuity give the Gaussian
with covariance $D_{R_0}$.
\end{proof}

\subsection{The Gaussian kernel and measurable cohomology}
\begin{theorem}[The kernel is an actual $L^2$ coboundary]
\label{thm:gaussian-kernel-coboundary}
For $R\in I$ and $v\in\R^4$, the following three statements are
equivalent:
\begin{enumerate}
\item $v^{\mathsf T}D_Rv=0$;
\item there is $b\in L^2(M,\nu)$ such that
 $v\cdot h_R=b-b\circ T_R$ almost everywhere;
\item there is $b_*\in L^2(Y_R^*,\nu_R^*)$ such that
 $v\cdot(G_R-\bar G_R)=b_*-b_*\circ F_R^*$ almost everywhere.
\end{enumerate}
If the first condition holds, one may choose $b$ with
$\|b\|_2\le C|v|$, uniformly in $R$, and take $b_*=b|_{Y_R^*}$.
\end{theorem}
\begin{proof}
Put $u=v\cdot h_R$ and $c_k=\int u(u\circ T_R^k)\dd\nu$.
By \eqref{eq:bv-correlation}, $|c_k|\le C|v|^2e^{-\gamma k}$.
If $v^{\mathsf T}D_Rv=0$, then $c_0+2\sum_{k\ge1}c_k=0$.
The exact variance identity gives
\[
 \int|S_{m,R}u|^2\dd\nu
       =-2\sum_{k=1}^{m-1}k c_k-2m\sum_{k\ge m}c_k
       \le C|v|^2 \qquad(m\ge1).
\]
Let $Uu=u\circ T_R$, a unitary operator on $L^2(\nu)$, and form
$b_M=M^{-1}\sum_{m=1}^M S_{m,R}u$. This is a bounded sequence in
$L^2$, with norm at most $C|v|$. Moreover
\[
 (I-U)b_M=u-\frac1M\sum_{m=1}^M U^m u\longrightarrow u
       \quad\hbox{in }L^2.
\]
Here the mean ergodic theorem applies; the collision map is ergodic
and $u$ has mean zero. A weakly convergent subsequence of $b_M$ has
limit $b$, and bounded linearity of $I-U$ shows $(I-U)b=u$.
This proves the second assertion and its norm estimate.

Discard a countable union of collision translates of the null set
where this identity or a representative fails. On the remaining
invariant full-measure set, telescope until the actual first return.
The compensation identity for one block gives
\[
 v\cdot(G_R-\bar G_R)(x)=b(x)-b(F_R^*x).
\]
The restriction belongs to $L^2(\nu_R^*)$ and has norm at most
$c_*^{-1/2}\|b\|_2$. Thus the second assertion implies the third.
Conversely, the third assertion telescopes for the stationary induced
map, so the $L^2(\nu_R^*)$ norm of
$v\cdot(J_{n,R}-n\bar G_R)/\sqrt n$ is at most
$2\|b_*\|_2/\sqrt n$. Its characteristic function tends to one.
Theorem~\ref{thm:actual-record-gaussian} identifies the same limit as
$\exp(-t^2v^{\mathsf T}D_Rv/2)$ for every real $t$, proving the first
assertion. This also completes the equivalence without presupposing
summability of induced correlations.
\end{proof}

\begin{remark}[What the low-frequency theorem supplies]
\label{rem:gaussian-versus-raw}
The matrix $D_R$ is now established by an absolutely convergent
collision Green--Kubo formula and is the covariance of the Gaussian
law of the true induced record. The zero-Gaussian-variance condition
has been identified with a genuine induced $L^2$ coboundary. A
measurable representative of its transfer function has not thereby
been made continuous at a selected periodic orbit. Thus the periodic
rank calculation cannot yet be applied to this representative to deduce
uniform positive definiteness. Neither convergence of the induced
Green--Kubo series nor convergence of its normalized second moments is
asserted by this stopping argument.

The estimate \eqref{eq:actual-record-gaussian} controls fixed compact
sets of rescaled central frequencies. It does not give the integrated
complementary-frequency bound in the raw local inversion theorem, an
induced eigenvalue expansion, or the complete critical/singular branch
sum. In particular it is not a raw density local limit, a weighted
local-density estimate, or a functional limit under a changed exact
conditioning event. These distinctions retain the original density
problem while supplying its actual Gaussian weak-limit input.
\end{remark}
